\section{Methodology}
\label{sec:methods}

The workflow is built entirely on open-source tools and sits within the
QGIS\,/\,\texttt{momepy} morphometrics ecosystem
\citep{fleischmann2019MomepyUrban, fleischmann2022EvolutioUrban} and the
\texttt{OSMnx} network ecosystem \citep{boeing2017Osmnx, boeing2025ModelingNetworks}.
Resilience indicators are produced with the open-source PlanX QGIS plugin (version 4.10) \citep{eminoglu2026planxresilience}, with spatial statistics run in
the PlanX suite's GeoStats module (version 0.9.17). Every algorithm parameter is
recorded, and the QGIS project, derived layers and code are released for
reproduction. The workflow joins a single commensurable analytical unit to a
two-stage explain-then-optimize synthesis (Figure~\ref{fig:flowchart}).

\subsection{Study area}
\label{sec:study-area}

The case is the \.{I}zmir functional urban region, a coastal Mediterranean
metropolis on the Aegean coast of western T\"urkiye. The region combines a dense
historic and planned core around the \.{I}zmir Gulf with apartment-block
expansion, hillside incremental settlement, waterfront transformation zones, and
an industrial-logistics and peri-urban edge. This cross-section of urban form is wide
enough to make a fabric-based analysis meaningful. It also carries a documented,
recent multi-hazard profile: extreme summer heat, pluvial flooding, and
coastal/sea-level processes, set against a seismically active context whose 2020
Aegean earthquake underlines the relevance of evacuation and assembly-area
accessibility. The 30 October 2020 Aegean earthquake (Mw~6.9) and
associated tsunami exposed critical shortcomings in evacuation routing and
assembly-area coverage, while severe pluvial flash floods struck low-lying
districts, particularly Bayrakl\i{} and Bornova, in consecutive years (2019,
2020, 2021), and summer urban-heat-island intensities routinely push
land-surface temperatures well above ambient
\citep{canguzel2024ClimateChange}.

The region's morphological range is operationalized into seven a-priori fabric strata:
historic core, grid residential, apartment-block fabric, waterfront transformation,
hillside/incremental tissue, industrial-logistics, and peripheral expansion.
These strata span a broad morphological and climatic gradient characteristic of a Mediterranean Csa (hot-summer) coastal metropolis. At one
extreme, the historic core exhibits the highest central density (FAR~1.69, 4.13~floors)
and near-maximal orientation entropy (0.91), yet its narrow-canyon geometry and
moderate ground coverage (0.37) keep mean summer land-surface temperature (LST) at
42.7\,$^{\circ}$C. The apartment-block fabric shows lower density (FAR~1.13,
3.49~floors, coverage~0.30) but records a higher LST of 43.7\,$^{\circ}$C, likely
reflecting larger impervious-surface fractions between blocks. Waterfront
transformation zones carry the tallest mean profile (5.08~floors) yet, benefiting
from coastal exposure (coverage~0.21), register the coolest stratum-mean LST of
41.7\,$^{\circ}$C. Hillside/incremental tissue and grid residential districts
occupy an intermediate range in both form (FAR from 0.34 to 0.80, 3.0 to 3.3~floors) and
thermal load (LST between 43.0 and 43.2\,$^{\circ}$C), while the industrial-logistics
stratum, despite the lowest building heights (1.88~floors), reaches the highest LST
of 45.5\,$^{\circ}$C owing to extensive impervious coverage and the coarsest,
least-vegetated grain. The peripheral expansion edge is the least dense (FAR~0.07,
1.87~floors) and the least built-up (coverage~0.03), with an LST of
42.4\,$^{\circ}$C driven primarily by bare-soil albedo. This
3.8\,$^{\circ}$C LST range across strata underscores why a fabric-resolved
analysis is essential for climate-sensitive planning in
\.{I}zmir \citep{canguzel2024ClimateChange}.

The spatial layout of the \.{I}zmir functional urban region and the geographical distribution of all 3{,}777 urban cells retained across these seven provisional fabric strata are mapped in Figure~\ref{fig:strata-map}.

\begin{figure}[H]
  \centering
  \includegraphics[width=\textwidth]{fig1}
  \caption{\.{I}zmir functional urban region study area map. (a) Context map showing district (ilçe) boundaries, coastline, and FUA boundary outline, with major central districts labeled. (b) Provisional fabric strata map showing the spatial layout of all 3{,}777 urban cells retained after exclusions (250\,m grid cells). Legend counts indicate the number of urban cells per stratum across the full census analysed in this study.}
  \label{fig:strata-map}
\end{figure}

\FloatBarrier
\subsection{Data sources}
\label{sec:data}
All input layers, their sources, resolutions, licences, role in the workflow and
current acquisition status are listed in the data inventory
(Table~\ref{tab:data-inventory}; see also Supplementary Figure~\ref{fig:supp-data-quality} for sensor resolutions and census distributions). Every input is open or openly licensed, which is
what allows the workflow, including its parameters and code, to be transferred to
other coastal metropolitan regions. The street network, building footprints, an
open digital elevation model and the regional boundary are assembled, reflecting the
growing availability of crowdsourced, open, and AI-derived datasets for building and
hazard mapping \citep{bhuyan2023MappingCharacte}; the satellite land-cover and
land-surface-temperature layers are obtained through a Google Earth Engine workflow,
enabling the hyperlocal mapping of surface urban heat islands using crowdsourced and
public satellite registers \citep{venter2020HyperlocMapping}, and the institutional
and demographic layers (green-blue infrastructure, assembly areas, essential services
and the T\"urkStat social-vulnerability inputs) are compiled from municipal open data
and national statistics.

\begin{table}[htbp]
  \centering
  \footnotesize
  \caption{Open-data and software inventory for the \.{I}zmir functional urban
  region. CRS for all spatial layers: EPSG:32635 (UTM Zone 35N).}
  \label{tab:data-inventory}
  \begin{tabular}{@{}p{2.7cm}p{3.6cm}p{2.4cm}p{1.6cm}p{3.4cm}@{}}
    \toprule
    Layer & Source & Resolution / date & Licence & Role in workflow \\
    \midrule
    Region boundary (FUR / districts)
      & National admin.\ boundaries; regional plan
      & Vector
      & Open
      & Study extent; stratification frame \\
    Analysis grid (250\,m; 500\,m test)
      & Generated (INSPIRE/GHSL-aligned)
      & 250 / 500\,m
      & Not applicable
      & Observation unit \\
    Street network
      & OpenStreetMap; \.{I}zmir Metropolitan Municipality (IMM) roads (207{,}734 segments)
      & Vector
      & ODbL / Open
      & Segment graph; network metrics; 400/800\,m network service areas \\
    Building footprints
      & OpenStreetMap (OSM); 901{,}609 building footprints with floor counts
      & Vector
      & ODbL
      & Morphological tessellation; built-form characters \\
    DEM / slope
      & Open elevation model (Copernicus DEM)
      & 30\,m
      & Open
      & Slope exclusion ($>15\%$); shade/SVF inputs \\
    Land cover / imperviousness
      & ESA WorldCover v200 (2021)
      & 10\,m
      & CC-BY 4.0
      & Impervious-exposure proxy; urban masking \\
    Land-surface temperature
      & Landsat 8/9 C2 L2 ST, summer (JJA) means 2014--2024
      & 100\,m $\rightarrow$ 250\,m
      & Open
      & \emph{Measured} dependent variable (SHAP explain stage) \\
    Coastline
      & OpenStreetMap (\texttt{natural=coastline})
      & Vector
      & ODbL
      & Coastal-exposure proxy \\
    Green and blue infrastructure
      & OpenStreetMap (landuse/leisure/natural)
      & Vector
      & ODbL
      & Cooling-access proxy \\
    Assembly areas, essential services
      & Municipal open data; OSM points of interest
      & Vector
      & Open / ODbL
      & Emergency- and daily-access proxies \\
    Social-vulnerability inputs
      & T\"urkStat ADNKS (31 Dec 2024): age structure, dependency ratio
      & Neighbourhood (\emph{mahalle})
      & Open
      & Social-vulnerability index; equity (Gini) \\
    \midrule
    Software
      & \multicolumn{4}{p{11.0cm}@{}}{QGIS 4.2; PlanX v4.10; PlanX GeoStats Lab v0.9.17; OSMnx;
      momepy; GeoPandas; NetworkX; XGBoost (gradient boosting); SHAP. Versions
      frozen for reproducibility.} \\
    \bottomrule
  \end{tabular}
\end{table}


\subsection{Analysis grid}
\label{sec:methods-grid}
The analytical surface is a 250\,m square grid (INSPIRE/GHSL-aligned) generated over the
\.{I}zmir functional urban region in the projected coordinate system EPSG:32635 (UTM Zone 35N).
The 250\,m cell resolution is chosen deliberately: although the Landsat 8/9 Thermal Infrared Sensor (TIRS) has a native spatial resolution of 100\,m (resampled to 30\,m by the USGS using cubic convolution), thermal data are aggregated to the 250\,m cells via zonal mean extraction to align with the scale of urban blocks and street networks. This cell size is large enough to yield stable morphometric and network statistics while remaining fine enough to resolve within-district spatial heterogeneity. Grid resolution is treated as an explicit object of analysis rather than a fixed assumption, and classifications are re-estimated at 500\,m to test scale stability (Section~\ref{sec:methods-cluster}).

To ensure the analysis describes inhabited urban fabric, grid cells that are predominantly non-urban, industrial-empty, water, or steeply sloped (mean slope $>15\%$) are excluded according to strict threshold rules. Specifically, cells with a built fraction below 0.10 (from ESA WorldCover 2021, 10\,m), a water fraction above 0.50, or a mean slope above 15\% (from the Copernicus GLO-30 DEM, 30\,m) are excluded. These filters reduce the total candidate grid from 16{,}496 cells to 3{,}777 retained urban cells.

All 3{,}777 retained urban cells enter downstream analysis as a \emph{full census} of the inhabited region without sampling. The seven a-priori fabric strata assigned to these cells are operationalized as rule-based proxies from land use, settlement history, and remote sensing prior to unsupervised clustering. They remain provisional and earmarked for refinement with planning documents and urban-design expertise; defining strata a priori and using clustering to test and refine them avoids the circularity of letting the classifier both define and confirm the typology (per-stratum cell counts are reported in Table~\ref{tab:fabric-profiles}).

\subsection{Cross-attribution}
\label{sec:methods-xattr}
The analytical observation is the sampled 250\,m grid cell. Three sources of
information are made commensurable within it. (i)~Building footprints are
partitioned into morphological tessellation cells, a cadastre-independent plot
proxy \citep{fleischmann2020Morphologic}, from which built-form characters are
computed and aggregated to the grid cell by area-weighted statistics (median and
interquartile range, plus densities). (ii)~Road centre-lines are cleaned and
planarized into a segment graph \citep{boeing2017Osmnx}, and network descriptors
are attributed to the cell they fall in or nearest to. (iii)~Accessibility and movement potential are evaluated within 400\,m and 800\,m
\emph{network service areas} around each cell. From the municipal road centre-lines we
build a routable street graph (nodes at segment endpoints, edges weighted by length)
and, for each cell, run a distance-capped Dijkstra search from the nearest node to
delimit the street segments reachable within 400\,m and 800\,m of network travel;
street-network configuration, including segment-length, node and intersection densities,
mean segment length and orientation entropy, is then measured over that service area.
Using network distance rather than a Euclidean buffer means the reach respects the
topology that governs walking access, so two cells equidistant ``as the crow flies''
but differently connected receive different accessibility profiles. This characterizes the street space surrounding each cell rather
than only what lies inside it, consistent with syntactic walkability and the
prediction of void street-interface occurrences \citep{tannous2021AccessibGreen,
de2023UnderstaPredicti}.
The cross-attribution rule maps each tessellation character $c$ onto grid cell
$g$ through the area-weighted aggregation
\begin{equation}\label{eq:xattr}
\bar{x}_{c}(g) \;=\; \frac{\sum_{t \in T(g)} a_{t \cap g}\, x_{c}(t)}
                        {\sum_{t \in T(g)} a_{t \cap g}},
\end{equation}
where $T(g)$ is the set of tessellation cells intersecting grid cell $g$,
$a_{t \cap g}$ is the area of the intersection between tessellation cell $t$ and
grid cell $g$, and $x_{c}(t)$ is the value of character $c$ on tessellation cell
$t$. The area-weighted median and interquartile-range analogues used for the
built-form characters are produced by the same operator, with weighted order
statistics in place of the weighted mean. Reach-based characters instead use the
network service area
\begin{equation}\label{eq:servicearea}
S_{r}(g) \;=\; \{\, e \in E \;:\; d_{\mathrm{net}}(v_{g}, e) \le r \,\},
\qquad r \in \{400, 800\}\,\mathrm{m},
\end{equation}
where $E$ is the set of street segments, $v_{g}$ is the graph node nearest to
the centre of cell $g$, and $d_{\mathrm{net}}$ is the edge-length-weighted
shortest-path (Dijkstra) distance along the street graph. The resulting feature
vector is the row on which all downstream analysis operates.
This rule is the methodological core: it is what lets movement, built form and
hazard be integrated in one unit rather than reported as parallel, unlinked
layers.

\label{sec:methods-indicators}
Network descriptors (computed with \texttt{OSMnx}/\texttt{NetworkX} and QGIS)
include connectivity and node density, meshedness, orientation entropy, circuity,
and centrality/configurational measures in the street-based and space-syntax
traditions \citep{porta2012StreetCentrali, araldi2019FromStreet,
boeing2022StreetNetwork, hillier1984SocialLogic}, read with awareness of the
known limits of axial/space-syntax abstraction at fine scales
\citep{pafka2020LimitsSpace}. Orientation entropy summarizes street-grid order as
the normalized Shannon entropy of segment bearings \citep{boeing2019UrbanSpatial},
\begin{equation}\label{eq:orient}
H_{o} \;=\; -\,\frac{1}{\ln 18}\,\sum_{i=1}^{18} p_{i} \ln p_{i},
\end{equation}
where $p_{i}$ is the length-weighted share of street segments falling in
orientation bin $i$ (18 bins; Table~\ref{tab:params}), so that $H_{o}=0$ for a
perfectly ordered single-direction grid and $H_{o}=1$ for a uniform distribution
of bearings. Morphometric descriptors (computed with
\texttt{momepy}/\texttt{GeoPandas}) include ground coverage, openness, compactness,
frontage continuity, block permeability, cell heterogeneity and open-space
fragmentation, following the numerical-taxonomy character set
\citep{fleischmann2021MeasurinUrban, fleischmann2022MethodolFoundati}. Each
indicator, its definition, the library function that computes it and its expected
direction are listed in the indicator table. The full indicator set comprises 29
variables (15 building/tessellation, 6 network, 5 contextual and 3 hazard
characters); their definitions, library functions and expected directions are
tabulated in the supplementary material.

\label{sec:methods-proxies}
We treat extreme heat, pluvial flooding and coastal exposure. Proxies sensitive
to urban morphology are mapped to each hazard: shade and solar-exposure
potential, sky-view factor, and the albedo of material scenarios along street
canyons for outdoor thermal comfort \citep{middel2018ViewFactor, mohammad2021EvaluatiRole};
green-blue cooling access and the role of water bodies in mitigating LST
\citep{wei2025NatureBased, cai2018WaterBodies, yu2020CriticalReview}; assembly-area
and daily-service access for the movement/coping axis; and the overlap of
impervious exposure with social vulnerability as an equity-sensitive exposure
proxy \citep{turner2025Guhvi, neumann2026ClimateJustice}. Exposure measures are
reported conservatively: impervious exposure is derived from satellite land cover rather
than from a building-density proxy, and built-form intensity is stated explicitly
where floor counts are unavailable. One boundary of scope is drawn here and
maintained throughout: the pluvial-flood layer (topographic susceptibility from
flow accumulation and depressions) is produced and released as a screening
overlay, but it does \emph{not} enter the prioritization stage. Entering a
topographic susceptibility surface into a ranking alongside measured heat would
give it a false equivalence with the better-constrained axes, so its promotion
to a need axis is deferred until it can be validated against hydrodynamic
modelling (Sections~\ref{sec:methods-optimize} and \ref{sec:disc-limits}).

\subsection{Explainable heat model}
\label{sec:methods-explain}
The dependent variable for the explain stage is \emph{measured} land-surface
temperature from summer Landsat thermal scenes, processed following established
Google Earth Engine surface-temperature workflows
\citep{ravanelli2018MonitoriImpact, zhou2018SatellitRemote}. At the grid-cell scale,
pooled across all 3{,}777 urban cells, a gradient-boosted-tree model (XGBoost) is
trained to predict land-surface temperature
from the morphometric and configurational drivers, with spatially blocked
cross-validation and reported $R^2$/RMSE against regularized linear baselines; a
pixel-scale attribution is deferred to future work. This builds on remote sensing studies
of LST variation, seasonal fluctuations, and the impacts of land architecture in
other contexts \citep{li2016RemoteSensing, haashemi2016SeasonalVariatio,
jiang2015AssessinImpacts, ranagalage2020SpatioteVariatio}. SHAP values then provide
global importance, dependence and interaction structure, and per-fabric local
attribution, yielding each fabric's \emph{dominant heat mechanism}. For feature
$j$, TreeSHAP defines the coalition-weighted marginal contribution
\citep{lundberg2020LocalExplanations} as
\begin{equation}\label{eq:shap}
\phi_{j} \;=\; \sum_{S \subseteq F \setminus \{j\}}
\frac{|S|!\,\bigl(|F|-|S|-1\bigr)!}{|F|!}\,
\Bigl[ f_{x}\bigl(S \cup \{j\}\bigr) - f_{x}(S) \Bigr],
\end{equation}
where $F$ is the feature set and $f_x(S)$ is the conditional expected model
output given coalition $S$. A
geographically weighted analysis of local non-stationarity is reserved for the
full study; a regularized linear baseline (ridge/elastic-net) guards against over-reading
collinear predictors. This follows the now-standard explainable-ML template for
the morphology$\rightarrow$temperature relationship, where the impacts of 2D/3D
morphology heterogeneity and density on LST have been extensively documented using
various grid- and block-based approaches \citep{wangZ2025PlainPlateau,
liu2026Nonlinear, li2026BlockGrid, guo2016CharacteImpact, guo2020ImpactUrban,
huang2019InvestigEffects, li2020InfluencDensity, jeon2023ImpactsUrban,
taleghani2015OutdoorThermal}. These methods also allow for testing the
climate impact of urban renewal, redevelopment, and high-resolution future climate
projections on thermal comfort and building energy performance \citep{qiao2020ImpactUrban,
renard2019EvaluatiEffect, hosseini2022HighResoluti, kabano2021EvidenceUrban}. This
workflow is applied at street/cell scale and fed forward into the typology.

\subsection{Typology and scale stability}
\label{sec:methods-cluster}
Indicators are direction-coded and z-standardized; principal component analysis
reduces collinearity; Ward hierarchical clustering then groups grid cells into
fabric-resilience profiles, with the number of clusters chosen by silhouette
analysis and dendrogram consistency. The resulting profiles are interpreted
against the a-priori strata and the wider fabric-typology literature
\citep{bobkova2021TowardsAnalytic, yu2023UrbanNeighbou, araldi2025MultiLevel,
debray2025UniversalPatterns, vartholomaios2025Detection}. Over the full census, the first
seven principal components explain 91\% of total variance; Ward linkage with the
cluster number chosen by silhouette analysis yields $k=4$ super-types (the
smallest retained cluster contains $n=123$ cells, so the minimum-cluster-size
guard of ten is not binding here) (Section~\ref{sec:results}).

\label{sec:methods-scale}
Because the choice of 250\,m is consequential, the classification is re-estimated
on a 500\,m grid and the two partitions are compared with the Adjusted Rand Index
and rank correlation; the 400/800\,m accessibility reach is optionally varied as a
secondary check. This reports which indicators are scale-sensitive and whether the
fabric-resilience typology is robust to resolution. Over the census the Adjusted
Rand Index is 0.38 (fair agreement): the broad fabric gradient persists across
resolution, though the 250\,m grid resolves a fourth, finer mode that the 500\,m grid
merges (Section~\ref{sec:results}).

\subsection{Spatial prioritization}
\label{sec:methods-spatial}
Spatial structure and equity are assessed in the PlanX suite's GeoStats module:
global
Moran's $I$ and local indicators of spatial association (LISA) on fabric clusters,
model residuals and priority scores \citep{sanchez2019SpotAnalysis,
purwanto2021SpatioteAnalysis}; Getis-Ord $G_i^{*}$ hot spots of adaptation
priority; and a Gini coefficient of the impervious-exposure $\times$
social-vulnerability burden as an explicit equity axis, tested for spatial
autocorrelation rather than decomposed. Multiscale geographically weighted
regression \citep{zhao2018GeographWeighted} belongs to the same suite and is
reserved for the local non-stationarity analysis of the full study; the
significance of the global and local statistics uses Monte-Carlo permutation.
With row-standardized $k$-nearest-neighbour weights, global spatial
autocorrelation is
\begin{equation}\label{eq:moran}
I \;=\; \frac{n}{\sum_{i}\sum_{j} w_{ij}}\;
\frac{\sum_{i}\sum_{j} w_{ij}\, z_{i} z_{j}}{\sum_{i} z_{i}^{2}},
\end{equation}
where $w_{ij}$ is a row-standardized $k$-nearest-neighbour weight ($k=8$)
and $z_i$ is the deviation of observation $i$ from the mean. The trapezoidal
Lorenz-curve approximation summarizes equity as
\begin{equation}\label{eq:gini}
G \;=\; 1 - \sum_{k=1}^{n} \bigl( X_{k} - X_{k-1} \bigr)
\bigl( Y_{k} + Y_{k-1} \bigr),
\end{equation}
where $X_k$ is the cumulative share of cells and $Y_k$ is the cumulative share
of the impervious-exposure $\times$ social-vulnerability burden.

\label{sec:methods-optimize}
Adaptation priority is framed as a multi-objective problem across the hazard and
capacity axes (heat, cooling deficit, access deficit, coastal exposure and social
vulnerability). The
heat axis is measured land-surface temperature; the SHAP attributions from
the explain stage are not entered into the optimizer as weights but are used to
interpret which morphological mechanism drives each priority fabric and
therefore which intervention type fits it (Section~\ref{sec:disc-policy}). The
Pareto-optimal frontier separates dominated fabrics from non-dominated trade-off
fabrics. With all axes coded so that larger values mean greater need, cell $i$
dominates cell $j$ when
\begin{equation}\label{eq:pareto}
i \succ j \;\iff\; n_{k}(i) \ge n_{k}(j)\;\; \forall k \in \{1,\dots,5\}
\;\;\wedge\;\; \exists\, k :\; n_{k}(i) > n_{k}(j),
\end{equation}
where $n_k(\cdot)$ denotes a need axis, and cells not dominated under this rule form
the priority frontier. Unambiguous intervention candidates are worse on every
axis;
because the fabric-level screen is degenerate when five axes meet only seven
fabrics, the dominance result is read at the cell level
(Section~\ref{sec:res-priority}). A
TOPSIS ranking \citep{hwang1981MultipleAttribute} is provided for decision use.
For the vector-normalized decision matrix, criterion entropy and its corresponding
weight are
\begin{equation}\label{eq:entropy}
e_{k} \;=\; -\frac{1}{\ln m} \sum_{i=1}^{m} p_{ik} \ln p_{ik},
\qquad
w_{k} \;=\; \frac{1 - e_{k}}{\sum_{l} \bigl(1 - e_{l}\bigr)},
\qquad
p_{ik} \;=\; \frac{r_{ik}}{\sum_{i} r_{ik}},
\end{equation}
and the same weighted matrix gives each cell's relative closeness to the ideal
solution,
\begin{equation}\label{eq:topsis}
C_{i} \;=\; \frac{D_{i}^{-}}{D_{i}^{+} + D_{i}^{-}},
\qquad
D_{i}^{\pm} \;=\; \Bigl( \sum_{k} \bigl( w_{k} r_{ik} - v_{k}^{\pm} \bigr)^{2} \Bigr)^{1/2},
\end{equation}
where $v_k^{+}$ and $v_k^{-}$ are the ideal and anti-ideal points. Rankings are
reported under equal-weight, entropy-weight and Monte-Carlo weight-perturbation
schemes, with 2{,}000 Dirichlet draws, rank-stability distributions and Kendall's
$\tau$. This multi-objective
framing follows established urban-heat and green-infrastructure optimization
\citep{zhang2024MultiObjective, zhu2025SpongeCities} and spatial prioritization
heuristics \citep{schroter2016SpatialPrioriti}, but is applied for the first time
at the level of urban fabric.

\label{sec:methods-repro}
Three controls accompany the workflow. First, the tessellation-as-plot-proxy
assumption can be validated against cadastral or plan records where these are
available. Second,
the hazard treatment is framed explicitly as screening (not hydrodynamic or
microclimate simulation), so outputs are read as relative priorities rather than
absolute risk. Third, all parameters are logged, all derived layers exported, and
the QGIS project and code released. All inputs are public or openly licensed
(OpenStreetMap, official municipal building and street layers, Copernicus/ESA land
cover, Landsat land-surface temperature and T\"urkStat demographics), so the workflow
carries no proprietary dependencies and transfers to other coastal metropolitan
regions with comparable open-data coverage.

\begin{figure}[H]
  \centering
  \includegraphics[width=\textwidth]{fig2}
  \caption{Methodological flowchart of the two-stage explain-then-optimize typomorphological urban resilience analysis pipeline, from raw inputs to multi-hazard priorities.}
  \label{fig:flowchart}
\end{figure}
